\subsection{A fixed logarithmic conductor bound for metabelian monodromy}

Throughout, $A$ is a complete normal two-dimensional local domain
of equal characteristic $p>0$ with algebraically closed residue field,
and $U$ is its punctured spectrum. Fix the cover $U_0\to U$ and
SNC resolution $X_0$ of its normalization supplied by F10b. Write
$D_0=\bigcup_i C_i$, $B=(C_i\cdot C_j)$,
\[
 k_i=\max\{0,2g(C_i)-2+\nu_i\},\qquad
 t=(-B)^{-1}k,\qquad T_0=\max_i t_i,
 \tag{F13.1}
\]
where $\nu_i$ is the number of nodes on $C_i$. As in F10, choose
a further blowup if necessary so that $D_0$ is nonempty.
The logarithmic breaks below are computed over the completed DVRs
at the generic points of $D_0$. The complete proof passed an
independent root audit; F13b also passed a separate BCE geometry audit.

\begin{lemma}[F13a: modular monomial representations with a specified normal subgroup]
Let $\Gamma$ be a finite group, let $N\triangleleft\Gamma$ be
abelian with $\Gamma/N$ abelian, and let $\kappa$ be an algebraically
closed field of characteristic $\ell$. Every nonzero simple
$\kappa[\Gamma]$-module is induced from a one-dimensional character
of a subgroup $K$ containing $N$. In particular $K$ is normal in
$\Gamma$ and $\Gamma/K$ is abelian. No restriction on the
$\ell$-part of $|\Gamma|$ is imposed.
\end{lemma}
\begin{proof}
The $\ell$-primary subgroup $N_\ell$ is normal in $\Gamma$ and
acts trivially by the normal-primary invariants argument of F12a,
with $p$ replaced by $\ell$. We may quotient it out. Restriction
to the resulting prime-to-$\ell$ abelian group $N$ is a direct
sum of character spaces. Simplicity makes the characters occurring
in $V$ a single $\Gamma$-orbit. Choose one character $\chi$, let
$I$ be its stabilizer, and put $W=V_\chi$. The character-space
sum gives $V=\operatorname{Ind}_I^\Gamma W$. Moreover $W$ is
simple as an $I$-module: a proper nonzero submodule would induce
a proper nonzero submodule of $V$.

Put $A_1=I/N$. The pairing
\[
 b:A_1\times A_1\longrightarrow\kappa^*,\qquad
 b(\bar x,\bar y)=\chi([x,y])
 \tag{F13.2}
\]
is well-defined, alternating and bimultiplicative. Indeed,
$[I,I]\subset N$, and the $I$-invariance of $\chi$ kills
$\chi([x,n])$ for $n\in N$; the commutator identities then
prove independence of lifts and bimultiplicativity. Its radical
contains the $\ell$-primary part of $A_1$, since $\kappa^*$ has
no nontrivial $\ell$-power torsion. Choose a maximal isotropic
subgroup $L\subset A_1$. It contains the radical and satisfies
$L=L^\perp$: an element of $L^\perp\setminus L$ could otherwise
be adjoined to produce a larger isotropic subgroup.

Let $K$ be the inverse image of $L$ in $I$. Commutators of
$K$ act on $W$ through the scalar $\chi([K,K])=1$, so the
finite family of matrices from $K$ commutes. Over $\kappa$ it
has a common eigenvector, defining a character
$\theta:K\to\kappa^*$ extending $\chi$. Frobenius reciprocity
gives a nonzero map
$\operatorname{Ind}_K^I\theta\to W$.
The stabilizer of $\theta$ in $I$ is $K$: for $x\in I$, equality
of $\theta$ and its conjugate means precisely
$b(\bar x,L)=1$, hence $\bar x\in L^\perp=L$.
Thus the restriction of $\operatorname{Ind}_K^I\theta$ to
$K$ is a direct sum of pairwise distinct character lines.
Every submodule is a sum of those lines, and $I$ acts transitively
on them. The induced module is therefore simple and the nonzero
map is an isomorphism. Induction in stages gives
$V=\operatorname{Ind}_K^\Gamma\theta$.
The preimage of $K$ before quotienting $N_\ell$ contains the
original $N$, as required. Finally every subgroup containing $N$
is normal in $\Gamma$, because $\Gamma/N$ is abelian.
\end{proof}

\begin{lemma}[F13b: the Swan bound is unchanged by a finite \'etale resolution cover]
Let $f:X_1\to X_0$ be any connected finite \'etale cover, and put
$U_1=f^{-1}(U_0)$ and $D_1=f^{-1}(D_0)$. If $C'_a$ lies over
$C_i$ with degree $d_a$, every finite $p$-primary character on
$U_1$ has Kato Swan conductor at $C'_a$ at most $t_i$.
In particular the bound is independent of $\deg f$.
\end{lemma}
\begin{proof}
The surface $X_1$ is regular and proper over $A_0$, the
normalization underlying $U_0$. Its Stein factor is the spectrum
of the normalization of $A_0$ in the connected cover $U_1$.
This normalization is finite, complete and local with residue
field $k$; $X_1$ is a resolution of it, and $D_1$ is its reduced
exceptional SNC divisor. These assertions can also be checked
from the finite \'etale morphism over $X_0$: away from $D_0$
it is exactly the specified finite \'etale cover, regularity
and the SNC condition are preserved, and the connected proper
closed fiber follows from proper henselian specialization.
Thus the proof of F10 applies to $X_1$.

Let $B'$ be the intersection matrix of the components $C'_a$.
The finite \'etale curve map $C'_a\to C_i$ gives
\[
 2g(C'_a)-2=d_a(2g(C_i)-2),\qquad \nu'_a=d_a\nu_i.
 \tag{F13.3}
\]
Consequently $k'_a=d_a k_i$. The projection formula gives, for
any rational vector $v=(v_i)$,
\[
 \left(B'f^*v\right)_a
 =\left(f^*\sum_i v_iC_i\right)\cdot C'_a
 =d_a(Bv)_i,
 \qquad (f^*v)_a=v_i.
 \tag{F13.4}
\]
Taking $v=t$ shows $(-B')f^*t=k'$, so
$(-B')^{-1}k'=f^*t$. Apply F10's refined-Swan inequality
$B's'\geq-k'$ and its nonnegative inverse-matrix argument:
$s'_a\leq(f^*t)_a=t_i$. This uses the same
Krishna--Majumder levelwise refined-Swan input as F10; it does
not assert a bound on the total order of the character.
\end{proof}

\begin{proposition}[F13: a uniform logarithmic conductor after one fixed cover; Proved]
For every prime $\ell\ne p$ and every finite locally constant
$\mathbf F_\ell$-sheaf $M$ on $U$ with metabelian monodromy,
the highest logarithmic break of $M|_{U_0}$ at $C_i$ is at most
$t_i$. The fixed cover and all the bounds are independent of
$\ell$, the rank and the finite metabelian monodromy group.
\end{proposition}
\begin{proof}
Let $\Gamma$ be the monodromy group of $M$, and let
$\Gamma_0$ be the image of $\pi_1(U_0)$. Put
$N=\Gamma_0\cap\Gamma'$. It is abelian and normal, and
$\Gamma_0/N$ is abelian. By F10b, the pullback of the original
abelian quotient $\Gamma^{\mathrm{ab}}$ extends over $X_0$.
Therefore the image of every boundary inertia group in
$\Gamma_0$ lies in this specified subgroup $N$.

Extend scalars to $\overline{\mathbf F}_\ell$ and take a
composition series of the restricted representation. For a
simple constituent $V$, F13a gives
$V=\operatorname{Ind}_K^{\Gamma_0}\theta$ with $K\supset N$.
The associated cover $U_1\to U_0$ is unramified at every
exceptional DVR, because all boundary inertia lies in $N$.
Purity on the regular surface $X_0$ extends its normalization
to a finite \'etale cover $f:X_1\to X_0$.
The character $\theta$ is a rank-one local system on $U_1$.
All its values lie in a finite extension of $\mathbf F_\ell$.
Split $\theta$ into its $p$-primary part and its part of order
prime to $p$. The latter is tame and has logarithmic break zero.
The former has Kato Swan conductor at $C'_a$ at most $t_i$
by F13b.

Kato--Saito, Theorem~1.3, identifies the rank-one Kato Swan
conductor with the Abbes--Saito logarithmic conductor over the
completed DVR, including imperfect residue fields
\cite[Theorem~1.3, p.~2]{KatoSaito2019}.
Identifying the finite cyclic character values with complex
roots of unity does not change its kernel or either conductor.
Thus the logarithmic break of $\theta$ at $C'_a$ is at most
$t_i$.

For precision, induction through this finite \'etale cover does
not multiply the break by its degree. At the completed DVR of
$C_i$, each local field extension belonging to a component of
$f$ is unramified. For any positive rational $r$, Abbes--Saito
base-change compatibility identifies its logarithmic group
$G^r_{\log}$ with the same subgroup of the original absolute
Galois group \cite[Proposition~3.15(3)]{AbbesSaito2002}.
For $r>t_i$, this subgroup acts trivially on each conjugate
character in the restricted induced representation, and trivially
on the cosets of every unramified local extension. It therefore
acts trivially on the induced representation. This proves the
claim for each simple constituent.

Finally the image of any positive logarithmic ramification group
is a finite $p$-group. Its representations in characteristic
$\ell\ne p$ are semisimple. Hence triviality on every constituent
implies triviality on the whole representation. Scalar extension
is faithful, so the same assertion holds for $M$. This last step
also shows that no semisimplicity of $M$ was assumed.
\end{proof}

\paragraph{The remaining complete-local Betti gate.}
F13 is a conductor estimate, not yet an $H^1$ estimate. The precise
additional assertion that would close the coefficient-$\ell\ne p$
case at derived length three is: for this fixed complete-local
pair $(X_0,D_0)$ and the fixed bounds $t_i$, there is a constant
$C$ such that
\[
 \dim_\Lambda H^1(U_0,L)\leq C\operatorname{rank}_\Lambda L
 \tag{F13.5}
\]
for all finite fields $\Lambda$ of characteristic different from
$p$ and all finite local systems $L$ whose logarithmic breaks
at $C_i$ are bounded by $t_i$. It would suffice to establish
(F13.5) just for the systems occurring in F13.
Inflation--restriction for the fixed $U_0/U$ and the complement
count of F11 would then give a uniform generator bound for all
quotients of derived length at most three. No such conclusion is
being made without (F13.5).

Hu--Teyssier prove the corresponding rank-linear estimate for a
proper smooth \emph{finite-type} $k$-scheme $\overline V$ and
$V=\overline V\setminus E$, uniformly in finite coefficient fields
\cite[Theorem~7.34, p.~41]{HuTeyssier2025Betti}. Their relative
Theorem~7.15 likewise requires finite-type schemes
(Definition~7.14 and Theorem~7.15, p.~33). The nearby-cycle
Theorem~4.7 requires a geometric henselian trait
(Definition~4.6 and Theorem~4.7, pp.~15--16), and its proof uses
that condition in the spreading-out step. Neither statement,
as printed, applies to the fixed complete-local pair above.
A comparison, approximation or local Betti theorem with constants
uniform for this pair remains necessary. This explicit scope
check prevents treating finite-type spreading out for one sheaf
as a uniform result for all ranks and all sheaves.
